\section{Introduction}\label{sec:introduction}

A proper $k$-coloring of the Euclidean plane is a function
$c:\R^2\to\{1,\ldots,k\}$ such that $c(x)\ne c(y)$ whenever
$\norm{x-y}=1$. The Hadwiger--Nelson problem asks for the least such
$k$, denoted by $\chi(\R^2)$, with no regularity imposed on the color
classes. Equivalently, it asks for the chromatic number of the graph
whose vertices are all points of the plane and whose edges are the
unit-distance pairs. Edges in a finite drawing of this graph may cross.

The problem dates to Edward Nelson in 1950. Nelson obtained the lower
bound four, and John Isbell found the upper bound seven by a hexagonal
coloring; the historical record of these unpublished observations is
recounted by \citet[pp.~15--19]{Soifer2003History}. Hadwiger's published
formulation \citep{Hadwiger1961} asks for a covering by sets that avoid
unit distance and gives the bounds four and seven. The seven-point
construction of \citet{MoserMoser1961}, now called the Moser spindle,
gives a particularly small finite obstruction to three colors. We use
its two overlapping-triangle constraints again in the final step of our
proof and verify the required placement explicitly.

The lower bound remained four until \citet{deGrey2018} constructed a
finite unit-distance graph that is not four-colorable in 2018.
His corrected construction has 1,581 vertices. Subsequent work made
these finite obstructions smaller and their verification more accessible:
\citet{Heule2018} used satisfiability proofs and extraction of
unsatisfiable cores to obtain 553-vertex examples, and
\citet{Parts2020Minimization} obtained a 509-vertex graph.
\citet{ExooIsmailescu2020} gave an alternative proof of the lower bound five,
and Parts gave a human-verifiable proof \citep{Parts2020Human}.
These advances concern finite obstructions to four colors. The
compactness theorem of \citet[Theorem~1]{deBruijnErdos1951} says that a
graph is $k$-colorable, for fixed finite $k$, if and only if every finite
subgraph is $k$-colorable. Consequently an unrestricted lower bound can
also be established without first displaying a finite graph witnessing
it.

The problem with measurable color classes has a different history.
Write $\chi_m(\R^2)$ for the least number of colors when each class must
be Lebesgue measurable. \citet{Falconer1981} proved
$\chi_m(\R^2)\ge5$ using density points and rotations.
\citet{Payne2009} developed this method for translation-invariant
unit-distance subgraphs. In particular, the graph on all of $\R^2$
whose edges have rational unit displacement vectors is two-colorable,
while every measurable coloring requires at least five colors. Thus
measurable and unrestricted colorings cannot be identified merely
because their distance constraints look the same.

For colorings by regular regions, \citet{Woodall1973} studied
distance-avoiding closed covers and region decompositions. Townsend
corrected Woodall's argument for a six-color lower bound in the
planar-map setting, announcing the repair in 1981 and giving its details
in 2005 \citep{Townsend1981,Townsend2005}. More recently,
\citet{SokolovVoronov2025} proved a lower bound seven for polygonal
colorings in a locally finite map framework. Their interface argument
excludes both adjacent colors wherever distances to the interface take
values both below and above one. We derive corresponding exclusions
from weak measurable data, without assuming regular color boundaries.

We prove the following result, working throughout in ZFC.

\begin{theorem}\label{thm:main}
The Euclidean plane has no proper five-coloring, even when arbitrary
color classes are allowed. Consequently,
\[
  6\le \chi(\R^2)\le 7.
\]
\end{theorem}

The remaining alternatives six and seven are unresolved.\footnote{
The May 14, 2026 manuscript of \citet{Reed2026} announces
$\chi(\R^2)=7$ through a circle-density argument. Its proposed strict
bound of $\pi/3$ on the angular measure of a unit-independent subset
of the unit circle fails for the half-open arc
$\{e^{it}:0\le t<\pi/3\}$: this arc has angular measure $\pi/3$
and contains no unit pair. The displayed formal theorem assumes that
density bound as a hypothesis, so the argument does not establish the
announced equality.}
The proof has two parts: a transfer from unrestricted proper colorings to a
measurable condition, and a geometric obstruction to five labels under that condition.
The transfer holds for every finite number of colors.

\begin{definition}[Weak measurable coloring]\label{def:weak-coloring}
Let $k$ be a positive integer, let $S^1\subset\R^2$ be the unit circle,
and let $\sigma$ be its normalized arc-length measure. A
\emph{weak measurable $k$-coloring} is a Lebesgue measurable function
$c:\R^2\to\{1,\ldots,k\}$, with color classes $A_i=c^{-1}(\{i\})$,
such that, for every $R>0$,
\[
  \sum_{i=1}^k\int_{B(0,R)}\int_{S^1}
  \ind_{A_i}(x)\ind_{A_i}(x+u)\,d\sigma(u)\,dx=0.
\]
The integrals use the completed product measure; equivalently, one
may first choose Borel representatives of the color classes forming
a partition of the plane.
\end{definition}

This definition permits exceptional unit pairs. It is unchanged by
modifying the coloring on a plane Lebesgue null set: by Fubini, the
pairs for which either $x$ or $x+u$ belongs to that set have product
measure zero. The required bridge to the unrestricted problem is the
following equivalence.

\begin{theorem}[Transfer of colorability]\label{thm:transfer}
For every positive integer $k$, in ZFC,
\begin{equation}\label{eq:source-1}
  \begin{gathered}
    \text{a proper $k$-coloring of $\R^2$ exists}
    \\ \Longleftrightarrow \\
    \text{a weak measurable $k$-coloring of $\R^2$ exists}.
  \end{gathered}
\end{equation}
\end{theorem}

\begin{theorem}\label{thm:weak-five}
There is no weak measurable five-coloring of the Euclidean plane.
\end{theorem}

\begin{proof}[Deduction of Theorem~\ref{thm:main}]
A proper five-coloring would give a weak measurable five-coloring by
Theorem~\ref{thm:transfer}, contrary to Theorem~\ref{thm:weak-five}.

For the upper bound, use the hexagonal construction described by
\citet{Hadwiger1961}.
Take Voronoi hexagons of a triangular lattice, each with circumradius
$r=2/5$. Their centers have nearest-neighbor distance $\sqrt3r$.
Write the lattice as a scaled copy of $\Z+\Z\omega$, where
$\omega=e^{2\pi i/3}$. Multiplication by $2-\omega$ gives a similar
sublattice of index seven, since $|2-\omega|^2=7$. Color centers by
its seven cosets, and give each point the color of its hexagon,
assigning boundary points to any incident hexagon. Distances inside
one hexagon are at most $2r<1$. Centers of distinct same-color
hexagons are at least $\sqrt{21}r$ apart, so points in those hexagons
are at least $(\sqrt{21}-2)r>1$ apart. This is a proper seven-coloring,
including all boundaries.
\end{proof}

\subsection{Proof overview and the role of earlier methods}

The two parts of the proof address different obstacles. A coloring with
arbitrary classes provides no density-point structure to use, while a measurable
coloring may have boundaries far too irregular to analyze as a planar
map. Theorem~\ref{thm:transfer} supplies the measurable input. The proof
of Theorem~\ref{thm:weak-five} then constructs the connected interfaces
it needs instead of assuming them. Figure~\ref{fig:proof-route} shows how
these two parts combine.

\paragraph{From arbitrary labels to measurable fields.}
Let $F$ be the field of real algebraic numbers, $E=F(i)$ the algebraic
plane, and $K=\{u\in E:|u|=1\}$ its group of algebraic rotations.
Restricting a proper coloring to the countable set $E$ and averaging
its translates and rotations gives an invariant probability law on
proper labelings of $E$. This is a standard amenable averaging
construction. Probabilistic approaches to the Hadwiger--Nelson problem
already use invariant randomization and finite-graph compactness
\citep{BourgeatHeinrichMelottiRobert2015,GwynStavrianos2020};
F\o lner averaging over countable groups of plane isometries also occurs
in the fractional-coloring work of
\citet{MatolcsiRuzsaVargaZsamboki2025}. These methods supply symmetry.
Here a further argument must retain the unit-distance exclusions while
producing ordinary measurable functions on the plane.

For this purpose let $D$ be the compact group of all characters of the
additive discrete group $E$. A character is a homomorphism from $E$ to
the unit complex circle. The ordinary continuous characters form the
Borel subgroup
\[
 C=\{z\mapsto\exp(i\xi\cdot z):\xi\in\R^2\}\subset D.
\]
Rotations act on characters by $(ud)(z)=d(uz)$.
The spectral theorem, in $L^2$ of the invariant labeling law,
separates each label indicator into the part
supported on $C$ and its orthogonal remainder. The central rigidity
statement, Theorem~\ref{thm:rigidity}, says that every $K$-invariant
probability on $D$ assigning zero mass to $C$ is Haar probability. The
remainder therefore has zero correlation under every nonzero algebraic
translation, so removing it preserves the zero unit correlations.
The $C$-supported subspace is the space of functions on a probability
factor, so projection onto it is conditional expectation and preserves
nonnegativity and keeps the sum of the label coordinates equal to one.
Translation operators on this factor extend continuously to $\R^2$.
The translated coordinates admit jointly measurable versions, giving
fields on the plane and the probability space. For a suitable sample,
choose the least label with positive coordinate at each point
where the coordinates sum to one, and assign arbitrary labels on the
remaining plane-null set. The zero unit correlations give the weak
measurable coloring condition. Section~\ref{sec:transfer} carries out this passage, including
the measurable-version and exceptional-pair details.

The rigidity proof occupies Sections~\ref{sec:spectral}
and~\ref{sec:rigidity}. Its structural ingredient is the
Furstenberg--Zimmer compact-extension method
\citep{Furstenberg1977,ZimmerExtensions1976,ZimmerDiscrete1976}.
We give the particular compact-factor tower and conditional averaging
arguments used here; the corrected modern account of
\citet{JamneshanStructure,JamneshanCorrigendum} supplies further
context. For an ergodic law, a relative singularity lemma along that
tower shows that the
law of the difference of two characters, sampled conditionally independently
over its terminal factor, still gives $C$ measure zero. Its Fourier
coefficient is nonnegative and radial, and its line averages defined
by an invariant mean vanish.
Two identities between sums of algebraic unit directions give a difference
inequality for this coefficient. A finite number-field argument shows
that a positive value at any nonzero radius would force a uniform positive
lower bound on an interval of algebraic radii, contradicting those averages.
This proves Haar rigidity without a regularity assumption on the original
coloring.

For the converse transfer, density points of a weak color class contain no unit pair
(Lemma~\ref{lem:palettes-density-proper}). This adapts the familiar
density-point mechanism of Falconer and Payne to product-almost-everywhere
unit-pair exclusion. A common
translation places any fixed finite configuration among points having
density one in their assigned color, called typical points,
and graph compactness then gives a proper coloring of the whole plane.
The weak condition uses Lebesgue measure in position and arc length in
direction; it is distinct from the invariant-mean notion of
monochromatic-pair frequency in \citet{GwynStavrianos2020}.

\paragraph{From weak labels to connected exclusions.}
For the five-color obstruction, fix a center $x$ and take weak-* limits
of the typical angular samples
$e\mapsto(\ind_{A_i}(y+re))_{i=1}^5$, where $e\in S^1$,
$y\to x$, and $r\to1$. The \emph{palette} at $(x,e)$ collects,
up to angular null sets, all labels having positive weight at $e$ in
these limits. Smooth transport of two unit directions gives exclusions
between palettes. Section~\ref{sec:palettes}
shows that almost every palette contains at most two labels, at every
fixed center. Section~\ref{sec:transitions} records abundant transitions
between labels in a finite graph at each center and proves that centers
where this graph contains a cycle form a closed locally finite set.
The estimates use the decay of the Fourier transform of circle measure;
no finite perimeter or smoothness of a color boundary is assumed.

Section~\ref{sec:interfaces} smooths the color indicators by disk
averages, then thresholds and renormalizes the resulting probability
vectors. Outside disks with vanishing total radii, these modified
vectors have at most two positive coordinates, so they lie in the
complete graph on the five labels. The forest
structure away from cyclic centers lets us fill the holes in those regions
continuously in that graph. A planar cover obstruction forces a
nontrivial loop around one cyclic-center disk. A simple cycle in the
reduced edge path of that loop supplies the label cycle. For each edge
midpoint, a connected component of its inverse image crosses the annulus.
Limits of these components give compact connected sets $K_j$ with a
common point $x$.
Each $K_j$ is assigned the two endpoint labels of its edge.
Both endpoint labels are absent almost everywhere from the open region
\[
 \Delta(K_j)=\{z:\min_{q\in K_j}|z-q|<1<\max_{q\in K_j}|z-q|\}.
\]
This connected-set exclusion has a geometric predecessor in the
interface arguments of \citet[Propositions~1--2]{SokolovVoronov2025}.
Here the continua and their common point are consequences of the
measure estimates and the topological extraction.

A simple cycle on five labels has length three, four, or five.
Section~\ref{sec:angular} excludes each possibility. Chosen limiting
directions of sequences in the continua approaching $x$ determine
allowed labels just inside and just outside the unit circle about $x$. A five-cycle gives
incompatible six-position angular words. A four-cycle gives two antipodal
pairs of directions and an impossible parity rule. A three-cycle forces six
alternating sectors for the two outside labels. The last configuration
leaves an open region using only the three cycle labels almost everywhere; a rational
placement certificate puts all seven vertices of the Moser obstruction
in that region. These contradictions complete the proof of
Theorem~\ref{thm:weak-five}.

\begin{figure}[htbp]
\centering
\begin{tikzpicture}[
  font=\small,
  stage/.style={draw=black!55, rounded corners=2pt, fill=black!3,
    align=center, text width=2.48cm, minimum height=1.05cm,
    inner sep=5pt},
  link/.style={-{Stealth[length=2mm]}, semithick, draw=black!70},
  node distance=0.48cm]
  \node[stage] (proper) {Arbitrary proper\\five-coloring};
  \node[stage, right=of proper] (law)
    {Invariant probability\\on proper labelings};
  \node[stage, right=of law] (factor)
    {Continuous\\translation factor};
  \node[stage, right=of factor] (weak)
    {Weak measurable\\five-coloring};
  \node[stage] (palettes) at ([yshift=-2.28cm]weak.center)
    {Angular palettes\\of size at most two};
  \node[stage] (cycle) at ([yshift=-2.28cm]factor.center)
    {Common exclusion\\continua on a cycle\\of length $3$, $4$, or $5$};
  \node[stage, text width=2.48cm,
    font=\footnotesize, minimum height=1.27cm] (cases)
    at ([yshift=-2.28cm]law.center)
    {$3$: Moser graph\\$4$: parity obstruction\\$5$: six-step words};
  \node[stage, fill=black!9] (end) at ([yshift=-2.28cm]proper.center)
    {Contradiction};
  \draw[link] (proper) -- (law);
  \draw[link] (law) -- (factor);
  \draw[link] (factor) -- (weak);
  \draw[link] (weak) -- (palettes);
  \draw[link] (palettes) -- (cycle);
  \draw[link] (cycle) -- (cases);
  \draw[link] (cases) -- (end);
\end{tikzpicture}

\caption{The contradiction argument for five colors. The transfer from
arbitrary proper colorings to weak measurable colorings is proved for
every finite number of colors; the geometric obstruction uses five labels.}
\label{fig:proof-route}
\end{figure}

\subsection{A probabilistic consequence}

\begin{samepage}
\begin{corollary}[Positive five-color badness]\label{cor:positive-badness}
There is a constant $\delta>0$ with the following properties.
\begin{enumerate}
\item[(i)] For every coloring $c:\R^2\to\{1,\ldots,5\}$ and every
normalized finitely additive left-invariant measure $\mu$ on all subsets
of the plane-isometry group $E(2)$, the invariant-mean badness satisfies
\[
 p_5^\mu(c):=\mu\bigl(\{T\in E(2):c(T^{-1}0)=c(T^{-1}1)\}\bigr)
 \ge\delta,
\]
where $\R^2$ is identified with $\C$.
\item[(ii)] If $c$ is a Lebesgue-measurable five-coloring with two
linearly independent periods, then
$\mathbb P(c(A)=c(A+U))\ge\delta$, where $A$ is uniform in a
fundamental parallelogram and $U$ is an independent uniform unit
direction. For any Lebesgue-measurable five-coloring, if the conditional
square-table frequencies
\[
 \mathbb P\bigl(c(A_R)=c(A_R+U)\mid A_R+U\in[-R,R]^2\bigr)
 \qquad(R>1)
\]
converge as $R\to\infty$, their limit is at least $\delta$; here
$A_R$ is uniform in $[-R,R]^2$ and independent of $U$.
\end{enumerate}
\end{corollary}
\end{samepage}

\begin{proof}
Theorem~\ref{thm:main} and graph compactness
\citep[Theorem~1]{deBruijnErdos1951} give a finite unit-distance graph
$H$ that is not five-colorable. Set $m=|E(H)|$ and $\delta=1/m$.
Every five-coloring of $H$ has at least one monochromatic edge.
Averaging this count with any of the measures in (i) gives
$mp_5^\mu(c)\ge1$, by isometry invariance and finite additivity;
this is the finite-graph bound of
\citet[Lemma~2.10 and Theorem~2.11]{GwynStavrianos2020}.
The periodic and asymptotic needle frequencies in (ii) are bounded
below by the minimum monochromatic-edge fraction of $H$, hence by
$1/m$, by \citet[Theorems~1--2]{BourgeatHeinrichMelottiRobert2015}.
Borel representatives give the same needle frequencies by Fubini;
for a periodic coloring choose them in one cell and extend periodically.
The graph and the positive constant here are existential.
\end{proof}
